\section{Cancellation between adjacent dimensions}\label{sec:cancellation}

We now choose the functions so that the adjacent-dimension terms in
$v$ almost cancel.  Product functions restricted to a simplex, with
the truncation controlled by moments of their squared density, also
occur in Maynard's sieve construction \cite[Section~7]{Maynard2015}.
Here we combine such functions at a short range of dimensions with
alternating signs.

\begin{lemma}[A one-dimensional profile]
\label{lem:small-first-moment}
For every $\lambda>0$ and $\beta>0$, there exists a nonnegative
$g\in C_c^\infty((0,\infty))$ such that
\[
 \int_0^\infty g(u)^2\,du=1,\qquad
 \int_0^\infty g(u)\,du=\sqrt\lambda,\qquad
 \mu_g:=\int_0^\infty u g(u)^2\,du<\beta.
\]
\end{lemma}

\begin{proof}
For $R\ge2$, choose a smooth $\phi_R$ on $(0,\infty)$, equal to $1/u$
on $[1,R]$, with $0\le\phi_R(u)\le1/u$ and support in $[1/2,2R]$.
Put
\[
 A_R=\int\phi_R,\qquad B_R=\int u\phi_R(u)^2\,du,
 \qquad N_R^2=\int\phi_R^2.
\]
Then $A_R\ge\log R$, $B_R\le\log(4R)$, and
$1/2\le N_R^2\le2$.  Set $g_0=\phi_R/N_R$,
$c=\lambda/(\int g_0)^2$, and $g(u)=c^{-1/2}g_0(u/c)$.
Changes of variable give exactly
\[
 \int g^2=1,\qquad \int g=\sqrt\lambda,
 \qquad
 \mu_g=\lambda\frac{B_R}{A_R^2}
       \le\lambda\frac{\log(4R)}{(\log R)^2}.
\]
Choose $R$ sufficiently large and then hold it fixed.
\end{proof}

The profile just constructed may have a very large support endpoint
and variance. We need only their finiteness for this one fixed
profile: the dimension will be chosen afterwards. The auxiliary
limit in the next proposition concerns the constants $w$ and $v$
attached to finite families, before any asymptotic in $X$ is taken.

\begin{proposition}[Cancellation between high dimensions]
\label{prop:dimension-cancellation}
For every fixed $\lambda>0$, there are families of the type used in
Proposition~\ref{prop:weight-moments}, with increasing finite maximum
dimension $k$, for which $w\to1$ and $v\to0$.
\end{proposition}

\begin{proof}
Fix $0<\beta<\gamma<\tau$ and choose $g$ from
Lemma~\ref{lem:small-first-moment}.  Write $d=\beta-\mu_g>0$.
Let $B$ be an upper endpoint of the support of $g$, and let $s_g^2$
be the variance of the probability density $g^2$.  All these
quantities are fixed and finite; they need not be bounded uniformly
in $\lambda$.

For a large integer $k$ and $0\le j\le k$, define
\[
 Q_j(\mathbf t)=k^{j/2}\prod_{i=1}^j g(kt_i),\qquad
 S_j=\sum_{i=1}^j t_i,
\]
with $Q_0=1$ and $S_0=0$.  Before imposing the simplex cutoff, these
products satisfy
\[
 \|Q_j\|_2=1,\qquad TQ_{j+1}=\sqrt{\lambda/k}\,Q_j\quad(0\le j<k).
\]
Since $\alpha_j/\alpha_{j+1}=(j+1)/\lambda$, alternating products
$(-1)^jQ_j/\sqrt{\alpha_j}$ at neighboring levels therefore cancel
up to the factor $1-\sqrt{(j+1)/k}$.  We use $r$ neighboring levels
near $k$: the condition $r/k\to0$ makes these factors small, while
$r\to\infty$ makes the two endpoint contributions small after
normalizing the squared mass equally across the levels.

Take $r=\lfloor\sqrt k\rfloor$ and $a=k-r+1$.  Choose once and for all a smooth cutoff
$\chi$ with $0\le\chi\le1$, equal to one on $[0,\beta]$ and zero on
$[\gamma,\infty)$.  At the $r$ indices $a\le j\le k$, set
\begin{equation}
 f_j(\mathbf t)=\frac{(-1)^j}{\sqrt r\sqrt{\alpha_j}}
                     Q_j(\mathbf t)\chi(S_j),
 \label{eq:active-levels}
\end{equation}
and set all other $f_j$, including $f_0$, to zero.  Each selected finite
family has compact interior support: every coordinate is bounded
away from zero by a positive support endpoint of $g$ divided by $k$,
and the cutoff forces $S_j\le\gamma<\tau$.

The density $Q_j^2$ has total mass one.  Under it the variables $kt_i$
are independent with density $g^2$, so, for $j\le k$,
\[
 \mathbb E_jS_j=\frac{j\mu_g}{k}\le\mu_g,
 \qquad \operatorname{Var}_j(S_j)=\frac{js_g^2}{k^2}
                                      \le\frac{s_g^2}{k}.
\]
Chebyshev's inequality therefore gives
\begin{equation}
 1-\frac{s_g^2}{kd^2}
 \le w=\frac1r\sum_{j=a}^k\mathbb E_j[\chi(S_j)^2]\le1.
 \label{eq:w-concentration}
\end{equation}
This proves $w\to1$.

For $a\le j\le k-1$, substituting $u=kt_{j+1}$ and using
$\alpha_j/\alpha_{j+1}=(j+1)/\lambda$ gives
\begin{align}
 \sqrt{\alpha_j}(f_j+Tf_{j+1})(\mathbf t)
 &=\frac{(-1)^jQ_j(\mathbf t)}{\sqrt r}
       \left[\chi(S_j)-\sqrt{\frac{j+1}{\lambda k}}
                  \int_0^\infty g(u)\chi(S_j+u/k)\,du\right].
 \label{eq:interior-cancellation}
\end{align}
For $k\ge2B/d$, Chebyshev's inequality gives, uniformly in $j\le k$,
\[
 \mathbb P_j(S_j>\beta-B/k)\le\frac{4s_g^2}{kd^2}.
\]
On the complementary event every cutoff in
\eqref{eq:interior-cancellation} equals one.  The bracket is then
$1-\sqrt{(j+1)/k}$, whose square is at most $(r/k)^2$.  Everywhere,
both terms in the bracket lie in $[0,1]$, since $g\ge0$,
$\int g=\sqrt\lambda$, and $j+1\le k$; its absolute value is at most
one.  After summing the $r-1$ interior terms, their total contribution
to $v/\lambda$ is at most
\[
 (r/k)^2+\frac{4s_g^2}{kd^2}.
\]

Here the probability estimate is with respect to $Q_j^2$, while the
coordinate-deletion integral contains $g(u)\,du$. The passage between
them uses no comparison of these measures: on the complementary event,
$S_j+u/k\le\beta$ holds for every $u$ in the support of $g$.
Thus the whole integral has its cutoff equal to one pointwise.

Exactly two further indices can contribute.  At the top, $j=k$, the
term is $\alpha_k\|f_k\|_2^2\le1/r$.  At the bottom,
$j=a-1=k-r$, only $Tf_a$ occurs, and
\[
 \sqrt{\alpha_{a-1}}Tf_a
 =\frac{(-1)^aQ_{a-1}}{\sqrt r}
      \sqrt{\frac a{\lambda k}}
      \int_0^\infty g(u)\chi(S_{a-1}+u/k)\,du.
\]
Its squared norm is at most $a/(kr)\le1/r$.  Consequently
\begin{equation}
 0\le\frac v\lambda
 \le\frac2r+\left(\frac rk\right)^2+\frac{4s_g^2}{kd^2}
 \longrightarrow0.
 \label{eq:v-cancellation-bound}
\end{equation}
Together with \eqref{eq:w-concentration}, this proves the proposition.
\end{proof}

\begin{proof}[Proof of Proposition~\ref{prop:weights}]
Fix $\lambda>0$ and one function $G$ as in
Proposition~\ref{prop:weight-moments}. Set
\[
 K_\lambda=C_*=\lambda+\lambda^2c_G^2.
\]
This choice is independent of $\varepsilon$. Given $\varepsilon>0$,
choose $\beta$, $\gamma$, $g$, and $\chi$ as in the preceding construction.
By Proposition~\ref{prop:dimension-cancellation}, one can select a
\emph{single finite} $k$ and its family with $w>0$ and
$v/w\le\varepsilon$.  Freeze this family and define
\[
 W_X(m)=\frac{Z(m)^2}{w}.
\]
All conclusions follow from Proposition~\ref{prop:weight-moments} by
division by $w$, or by $w^2$ for the second moment.  Once the construction
is chosen in terms of $\lambda$ and $\varepsilon$, the fixed-family
second-moment bound may be written $O_{\lambda,\varepsilon}(1)$.

The parameter order is essential: the auxiliary limit $k\to\infty$
produces one family, and only then does $X\to\infty$.  No divisor-sum
asymptotic uniform in $k$ is required.  For an arbitrary fixed gap
threshold $C>0$, the application in Section~\ref{sec:counting} chooses a fixed
$\lambda>\max\{1,C\}$ before making these choices.
\end{proof}
